%%%%%%%%%%%%%%%%%%%%%%%%%%%%%%%%%%%%%%%%%%
\section{Conclusions}\label{sec:conclusions}

\rev{We evaluated a coefficient-space PINN for the Allen--Cahn equation, in which a network of time alone outputs the coefficients of a Chebyshev expansion, against a certified reference solution, a space--time PINN and the classical Chebyshev collocation solver, on three benchmarks and with five seeds. The following statements are supported by the experiments.}

\rev{\emph{Accuracy.} With the exact Neumann basis the network reaches $L^2=\N{b1.fh_K32.l2}$ on the reaction-dominated benchmark B1, a factor \N{b1.ratio.van_fhK32} better than the vanilla PINN, but a factor \N{b1.ratio.fh_K32_cheb80} worse than classical collocation, which also needs three orders of magnitude less time; on B1 a model that neglects diffusion is already as accurate as the network to within a factor \N{b1.ratio.ode_fh_K32}. \TODO{one sentence on B2 and B3 once the runs exist.} The method is not a replacement for classical solvers on the forward problem.}

\rev{\emph{Structure.} The Neumann condition and, in two dimensions, the reflection symmetry can be imposed exactly through the basis; the two-dimensional symmetry violation stays at round-off ($\le\N{s2d.smax}$) in all configurations, while the vanilla PINN's grows monotonically. The energy of the learned trajectory obeys an exact a-posteriori balance (machine-checked in Lean~4): it can increase only through the PDE residual and the boundary defect, whose time integral is \N{bal.int_r_hard} for the revised $K=32$ network, a fraction \N{bal.frac_hard} of the energy decrease. The energy penalty has no measurable effect when the residual is small. Energy decay does not guarantee admissibility: the networks overshoot $|u|\le1$ by up to \N{mbp.excess_all}.}

\rev{\emph{Corrections of the first version.} The reported degradation of accuracy with the number of modes was caused by a boundary penalty whose effective weight decayed like $K^{-4}$, and disappears with an exact Neumann basis (or a correctly scaled penalty); the reported $38\%$ bias of the recovered interface parameter was the footprint of the starting value of the optimization, and an estimator based on variable projection recovers $\varepsilon$ to \N{inv.vp_err_005} and \N{inv.vp_err_01} for noise-free data at $\varepsilon=0.05$ and $0.1$, while for $\varepsilon=0.01$ the problem is not identifiable at the noise levels considered. The initial condition of the first version was incompatible with the boundary condition; all experiments were repeated with a compatible one.}

\rev{\emph{Open questions.} The interface regime at $\varepsilon=0.01$, which needs $K\ge128$ and long training, and settings in which a coefficient network could offer something that a classical solver cannot (parametric families, partially unknown physics, long-horizon prediction) were not tested. A proof of convergence of the trained network would require, besides the a-posteriori balance, a stability estimate of the PDE with respect to the residual.}
